\documentclass[11pt]{article}
\usepackage[margin=1in]{geometry}
\usepackage{amsmath,amssymb,amsthm}
\usepackage{graphicx}
\usepackage[hidelinks]{hyperref}
\setlength{\emergencystretch}{2em}

\newtheorem{theorem}{Theorem}
\newtheorem{lemma}{Lemma}
\newtheorem{corollary}{Corollary}
\newtheorem{remark}{Remark}

\newcommand{\C}{\mathbb C}
\newcommand{\dist}{\operatorname{dist}}
\newcommand{\Prim}{\operatorname{Prim}}

\title{Countable Capture for Von Neumann Algebras}
\author{Run \texttt{fa\_banach\_001}, lane 19}
\date{June 17, 2026}

\begin{document}
\maketitle

\section*{Source Problem}

The source paper is Jinghao Huang, Karimbergen Kudaybergenov, and Rui Liu,
\emph{The ball-covering property of von Neumann algebras and noncommutative
symmetric spaces}, arXiv:2311.14257.  It proves that a von Neumann algebra
has BCP, equivalently UBCP, if and only if it is atomic and can be represented
on a separable Hilbert space.  The same paper asks for a broader
noncommutative \(C^*\)-algebra version:

\begin{center}
\includegraphics[width=0.92\textwidth]{figures/open_problem_crop.png}
\end{center}

\section*{Statement}

\begin{theorem}
\label{thm:main}
For a von Neumann algebra \(M\), the following are equivalent:
\begin{enumerate}
    \item \(M\) has the countable-capture property;
    \item \(M\) fails BCP;
    \item \(M\) is not atomic and separably represented.
\end{enumerate}
\end{theorem}

Here countable capture is understood in the Banach-space sense: for every
countable \(C\subset M\), there is \(u\in S_M\) such that
\[
        \dist(c,\C u)=\|c\|
        \qquad(c\in C).
\]
This is equivalent to the contraction formulation in the packet
\[
\texttt{2311.14257\_countable\_capture\_characterization}.
\]

\section*{Idea of the Proof}

Countable capture always implies failure of BCP.  The source paper gives the
von Neumann BCP classification, so it remains to prove that every von Neumann
algebra outside the atomic separably represented class has countable capture.

There are two mechanisms.  In the atomless part, the source singular-state
lemma gives, for each centre \(x_n\), a singular state \(\varphi_n\) with
\[
        \varphi_n(x_n^*x_n)=\|x_n\|^2.
\]
A standard projection characterization of singular functionals lets us find
one nonzero projection \(p\) with \(\varphi_n(p)=0\) for all \(n\).  Then the
functionals
\[
        y\mapsto \frac{\varphi_n(x_n^*y)}{\|x_n\|}
\]
norm \(x_n\) and all vanish at \(p\), giving capture by the common-kernel
form of the characterization.

In the atomic nonseparable case, either there are uncountably many atomic
summands, which gives primitive-topological capture, or there is a summand
\(B(H)\) with \(H\) nonseparable.  A countable family of operators on \(H\)
has a separable reducing subspace that norms all of them; a projection
orthogonal to that subspace gives the common capture direction.

\section*{Singular States and Common Kernels}

We use the following standard characterization of singular positive
functionals on von Neumann algebras.

\begin{lemma}[common projection for countably many singular states]
\label{lem:singular-common}
Let \(N\) be a nonzero von Neumann algebra and let
\((\varphi_n)\) be countably many singular states on \(N\).  Then there is a
nonzero projection \(p\in N\) such that
\[
        \varphi_n(p)=0\qquad(n\ge1).
\]
\end{lemma}

\begin{proof}
Fix a normal state \(\psi\) on \(N\).  A singular positive functional
\(\varphi\) admits an increasing net of projections \(e_i\uparrow 1\) with
\(\varphi(e_i)=0\).  Hence, by normality of \(\psi\), for every
\(\varepsilon>0\) there is a projection \(e\) such that
\[
        \varphi(e)=0,\qquad \psi(1-e)<\varepsilon.
\]
Apply this to \(\varphi_n\) with \(\varepsilon_n=2^{-n-2}\), obtaining
projections \(e_n\) with \(\varphi_n(e_n)=0\) and
\(\psi(1-e_n)<2^{-n-2}\).  Put
\[
        p=\bigwedge_{n=1}^\infty e_n.
\]
Then
\[
        \psi(1-p)
        =
        \psi\Big(\bigvee_n(1-e_n)\Big)
        \le
        \sum_n\psi(1-e_n)
        <
        \frac12.
\]
Thus \(p\ne0\).  Since \(p\le e_n\), we have
\(\varphi_n(p)\le\varphi_n(e_n)=0\) for all \(n\).
\end{proof}

\begin{lemma}[atomless summands have capture]
\label{lem:atomless}
Let \(N\) be an atomless von Neumann algebra.  Then \(N\) has countable
capture.
\end{lemma}

\begin{proof}
Let \(C=\{x_1,x_2,\ldots\}\subset N\) be countable.  Discard the zero
elements.  By the singular-state norming lemma from the source paper, for
each \(x_n\ne0\) there is a singular state \(\varphi_n\) on \(N\) such that
\[
        \varphi_n(x_n^*x_n)=\|x_n\|^2.
\]
By Lemma~\ref{lem:singular-common}, choose a nonzero projection \(p\in N\)
with \(\varphi_n(p)=0\) for every \(n\).  Replacing \(p\) by its support
projection if necessary is unnecessary; already \(\|p\|=1\).

For each \(n\), define
\[
        f_n(y)=\frac{\varphi_n(x_n^*y)}{\|x_n\|},
        \qquad y\in N.
\]
By Cauchy--Schwarz for the positive functional \(\varphi_n\),
\[
        |f_n(y)|
        \le
        \frac{\varphi_n(x_n^*x_n)^{1/2}\varphi_n(y^*y)^{1/2}}{\|x_n\|}
        \le
        \|y\|.
\]
Moreover
\[
        f_n(x_n)=\|x_n\|,
\]
so \(f_n\) is a norming functional for \(x_n\).  Finally,
\[
        |f_n(p)|^2
        =
        \frac{|\varphi_n(x_n^*p)|^2}{\|x_n\|^2}
        \le
        \frac{\varphi_n(x_n^*x_n)\varphi_n(p)}{\|x_n\|^2}
        =
        0.
\]
Thus all \(f_n\)'s vanish at the same nonzero vector \(p\).  The common-kernel
form of countable capture applies.
\end{proof}

\section*{Atomic Nonseparable Cases}

\begin{lemma}
\label{lem:nonsep-bh}
If \(H\) is a nonseparable Hilbert space, then \(B(H)\) has countable
capture.
\end{lemma}

\begin{proof}
Let \(C=\{T_1,T_2,\ldots\}\subset B(H)\) be countable.  Choose a separable
closed subspace \(L\subset H\) which reduces every \(T_n\) and satisfies
\[
        \|T_n|_L\|=\|T_n\|\qquad(n\ge1).
\]
This is obtained by starting with countably many almost norming vectors for
the \(T_n\)'s and closing under all \(T_n\) and \(T_n^*\).  Since \(H\) is
nonseparable, \(L^\perp\ne0\).  Let \(p\) be a nonzero projection with
range contained in \(L^\perp\).

For every scalar \(\lambda\),
\[
        (T_n-\lambda p)|_L=T_n|_L.
\]
Hence
\[
        \|T_n-\lambda p\|
        \ge
        \|T_n|_L\|
        =
        \|T_n\|.
\]
Taking the infimum over \(\lambda\) gives
\[
        \dist(T_n,\C p)=\|T_n\|,
\]
so \(B(H)\) has countable capture.
\end{proof}

\begin{lemma}
\label{lem:atomic}
Let \(M\) be an atomic von Neumann algebra which is not separably represented.
Then \(M\) has countable capture.
\end{lemma}

\begin{proof}
Write
\[
        M\cong \prod_{i\in I} B(H_i)
\]
as an atomic von Neumann algebra.

If \(I\) is uncountable, then the coordinate ideals give an uncountable family
of pairwise orthogonal nonzero closed two-sided ideals.  Hence \(\Prim(M)\)
has no countable \(\pi\)-base, and the primitive-space countable-capture
theorem from the packet
\[
\texttt{2311.14257\_countable\_capture\_characterization}
\]
implies countable capture.

If \(I\) is countable, nonseparable representability means that at least one
summand \(B(H_{i_0})\) has \(H_{i_0}\) nonseparable.  Let
\[
        C=\{x_1,x_2,\ldots\}\subset M
\]
be countable, and write \(x_n=(x_{n,i})_{i\in I}\).  Apply
Lemma~\ref{lem:nonsep-bh} in \(B(H_{i_0})\) to the countable family
\((x_{n,i_0})\), obtaining a projection \(p\in B(H_{i_0})\) such that
\[
        \dist(x_{n,i_0},\C p)=\|x_{n,i_0}\|
        \qquad(n\ge1).
\]
Let \(u\in M\) be \(p\) in the \(i_0\)-coordinate and \(0\) elsewhere.  Then
\(\|u\|=1\).  For every scalar \(\lambda\),
\[
        \|x_n-\lambda u\|
        =
        \max\Big\{
            \sup_{i\ne i_0}\|x_{n,i}\|,
            \|x_{n,i_0}-\lambda p\|
        \Big\}
        \ge
        \max\Big\{
            \sup_{i\ne i_0}\|x_{n,i}\|,
            \|x_{n,i_0}\|
        \Big\}
        =
        \|x_n\|.
\]
Therefore \(\dist(x_n,\C u)=\|x_n\|\) for all \(n\), giving countable
capture.
\end{proof}

\section*{Proof of the Classification}

\begin{proof}[Proof of Theorem~\ref{thm:main}]
If \(M\) has countable capture, then \(M\) fails BCP by the general
countable-capture obstruction packet.

The source paper proves that a von Neumann algebra has BCP if and only if it
is atomic and separably represented.  Thus BCP failure is equivalent to being
outside that class.

It remains to show that every von Neumann algebra outside that class has
countable capture.  Let \(z\) be the central projection supporting the
atomless part of \(M\).  If \(z\ne0\), then \(zM\) has countable capture by
Lemma~\ref{lem:atomless}.  Given a countable \(C\subset M\), apply that
lemma to \(zC\) and obtain \(u\in zM\) with \(\|u\|=1\) and
\[
        \dist(zc,\C u)=\|zc\|\qquad(c\in C).
\]
Since \(u\) is zero on \((1-z)M\), for every scalar \(\lambda\),
\[
        \|c-\lambda u\|
        =
        \max\{\|(1-z)c\|,\ \|zc-\lambda u\|\}
        \ge
        \max\{\|(1-z)c\|,\ \|zc\|\}
        =
        \|c\|.
\]
Hence \(M\) has countable capture.

If \(z=0\), then \(M\) is atomic.  Since \(M\) is outside the atomic
separably represented class, it is atomic but not separably represented.
Lemma~\ref{lem:atomic} gives countable capture.
\end{proof}

\section*{Scope and Use for the C*-Problem}

This gives a complete countable-capture classification for von Neumann
algebras:
\[
        \text{BCP failure}
        \quad\Longleftrightarrow\quad
        \text{countable capture}.
\]
It supports the broader program of finding whether C*-algebraic BCP failure
is always witnessed by countable capture, while also showing the different
mechanisms needed: primitive topology, nonseparable atomic Hilbert summands,
and singular-state capture in atomless summands.

\section*{Novelty and Search Bounds}

The source paper already proves the von Neumann BCP/UBCP classification and
the singular-state norming lemma for atomless von Neumann algebras.  The local
indexes were searched for countable capture, von Neumann algebras, singular
states, atomless, atomic, and separably represented.  Web searches on June 17,
2026 for BCP and von Neumann classification found the source paper but no
separate countable-capture reformulation.

\section*{Human Review Focus}

Reviewers should check:
\begin{enumerate}
    \item Lemma~\ref{lem:singular-common}, especially the standard singular
          functional projection characterization;
    \item the use of the source singular-state norming lemma in
          Lemma~\ref{lem:atomless};
    \item the atomic nonseparable decomposition in Lemma~\ref{lem:atomic};
    \item the final central-summand maximum-norm argument.
\end{enumerate}

\begin{thebibliography}{9}

\bibitem{HKL}
Jinghao Huang, Karimbergen Kudaybergenov, and Rui Liu,
\emph{The ball-covering property of von Neumann algebras and noncommutative
symmetric spaces},
arXiv:2311.14257, 2023.

\bibitem{TakesakiI}
Masamichi Takesaki,
\emph{Theory of Operator Algebras I},
Springer, 1979.

\bibitem{lane19capture}
Run \texttt{fa\_banach\_001}, lane 19,
\emph{When the Countable-Capture BCP Obstruction Exists},
solution packet
\texttt{2311.14257\_countable\_capture\_characterization}, 2026.

\bibitem{lane19primitive}
Run \texttt{fa\_banach\_001}, lane 19,
\emph{A Primitive-Ideal Pi-Base Obstruction for \(C^*\)-Algebras with BCP},
solution packet
\texttt{2311.14257\_primitive\_pi\_base\_obstruction}, 2026.

\end{thebibliography}

\end{document}
